\documentclass[12pt,letterpaper]{article}

\usepackage[spanish]{babel}
\usepackage[utf8]{inputenc}
\usepackage[T1]{fontenc}
\usepackage{geometry}
\usepackage{longtable}
\usepackage{array}
\usepackage{hyperref}
\usepackage{graphicx}
\usepackage{xcolor}

\geometry{margin=2.5cm}

\title{Informe Técnico de Avance y Aseguramiento de Calidad\\ Plataforma Digital ILP Backend}
\author{Proyecto ILP Backend}
\date{2026}

\begin{document}
\maketitle
\tableofcontents
\newpage

\section{Resumen Ejecutivo}
El proyecto ILP Backend corresponde a una plataforma digital educativa basada en arquitectura modular y microservicios. Actualmente el build global del proyecto compila exitosamente con Gradle, lo cual permite avanzar hacia una fase de aseguramiento de calidad, arquitectura sostenible y DevOps.

\section{Estado Actual del Proyecto}
El backend contiene módulos como auth-service, gateway-service, discovery-server, student-service, tenant-service, assessment-service, adaptive-education-service, report-service, notification-service, monitoring-service, common-events y common-libs.

\section{Marco de Calidad}
El aseguramiento de calidad se estructura tomando como referencia IEEE 730, integrando planificación, control, verificación, validación, auditoría, trazabilidad, gestión de defectos y mejora continua.

\section{Estrategia de Pruebas}
La estrategia se basa en la pirámide de Cohn:
\begin{itemize}
  \item Pruebas unitarias: rápidas, aisladas, patrón AAA.
  \item Pruebas de integración: validan persistencia, servicios y contratos.
  \item Pruebas end-to-end: flujos críticos con Screenplay.
\end{itemize}

\section{Arquitectura}
Se propone evolucionar gradualmente hacia arquitectura hexagonal, bajo acoplamiento, alta cohesión y separación clara entre dominio, aplicación, infraestructura y presentación.

\section{DevOps e IaC}
Se incorporarán pipelines, análisis estático, cobertura, quality gates, contenedores, infraestructura como código y automatización reproducible.

\section{Gestión de Defectos}
Todo bug debe registrarse con identificador, severidad, prioridad, módulo, pasos para reproducir, evidencia, causa raíz, estado y responsable.

\section{Conclusiones}
El proyecto cuenta con una base funcional. La siguiente fase debe enfocarse en calidad medible, trazabilidad, pruebas automatizadas, análisis estático y arquitectura sostenible.

\end{document}
